\documentclass[12pt]{article}
\usepackage[utf8]{inputenc}
\usepackage[T1]{fontenc}
\usepackage[margin=1in]{geometry}
\usepackage{setspace}
\usepackage{xcolor}
\usepackage[colorlinks=true, urlcolor=blue!60!black, linkcolor=blue!60!black]{hyperref}
\setlength{\parindent}{2em}
\setlength{\parskip}{1.25em}
\renewcommand{\baselinestretch}{1.0}
\begin{document}
\begin{center}
{\LARGE \textit{Snyder v. Phelps}, 562 U.S. 443 (2011)\par}
\vspace{0.8em}
{\large Case Brief --- Four-Page Edition\par}
\end{center}
\vspace{1em}

\section{Memory Jogger}

Westboro Baptist Church picketed, on public land 1,000 feet away, the funeral of a Marine killed in Iraq; an 8--1 Court held that the First Amendment barred tort liability for peaceful picketing on matters of public concern on public land, considering the signs' content, form, and context and the protest's distance from the service (\href{https://supreme.justia.com/cases/federal/us/562/443/}{562 U.S. 443}).

\section{Detailed Case Facts}

Fred Phelps founded the Westboro Baptist Church in Topeka, Kansas, in 1955; its members picketed nearly 600 funerals to publicize their belief that God punishes the United States for tolerating homosexuality, especially in the military, and to condemn the Catholic Church. When Marine Lance Corporal Matthew Snyder was killed in Iraq, his father, Albert Snyder, arranged a Catholic funeral in Westminster, Maryland.

Phelps and six relatives notified the police and, following their directions, stood behind a temporary fence on a 10-by-25-foot plot of public land about 1,000 feet from the church. For about 30 minutes before the service they sang hymns and displayed signs such as ``Thank God for Dead Soldiers,'' ``Priests Rape Boys,'' ``God Hates Fags,'' and ``You're Going to Hell.'' No one entered church property, yelled, or used violence. The procession passed within 200 to 300 feet of the site; Snyder saw only the tops of the signs on the way to the service and read them on the news that night. Weeks later Westboro posted an online ``epic'' denouncing the Snyders, which the Court declined to consider because Snyder had not raised it in his petition.

\section{Procedural History}

Snyder sued Phelps, two of his daughters, and the church in federal court in Maryland. The district court granted summary judgment to Westboro on defamation and publicity given to private life, but a jury found for Snyder on intentional infliction of emotional distress (IIED), intrusion upon seclusion, and civil conspiracy, awarding \$2.9 million in compensatory and \$8 million in punitive damages, later reduced to a \$5 million total. The Fourth Circuit reversed, finding the picketing protected. The Supreme Court granted certiorari, heard argument on October 6, 2010, and decided the case on March 2, 2011.

\section{Judicial Votes}

Affirmed, \textbf{8--1}.
\begin{itemize}
\item \textbf{Majority:} Chief Justice Roberts, joined by Justices Scalia, Kennedy, Thomas, Ginsburg, Breyer, Sotomayor, and Kagan.
\item \textbf{Concurrence:} Justice Breyer.
\item \textbf{Dissent:} Justice Alito.
\end{itemize}

\section{Holding}

The First Amendment shields Westboro from liability for IIED, intrusion upon seclusion, and civil conspiracy arising from peaceful picketing on matters of public concern in a traditional public forum. The captive-audience doctrine does not permit recovery where mourners could see only the tops of distant signs and the service was not disrupted.

\section{Analysis of Opinions}

\textbf{Majority (Roberts).} Under \textit{Hustler v. Falwell}, \textit{Connick}, and \textit{Dun \& Bradstreet}, speech on public issues occupies the ``highest rung of the hierarchy of First Amendment values,'' judged by its ``content, form, and context.'' Westboro's signs, though crude, addressed the nation's moral conduct, homosexuality in the military, and clergy scandals, and their ``overall thrust and dominant theme'' was public. The funeral context did not make the speech private, because Westboro had no prior grudge against the Snyders and had picketed on these themes for 20 years. It stood on public land next to a public street, ``the archetype of a traditional public forum.'' Content-neutral buffer laws raise separate questions, but Maryland's ``outrageousness'' standard is too subjective and invites juries to punish unpopular views. The captive-audience doctrine applies sparingly, mainly in the home, and the conspiracy claim fell with the others.

\textbf{Breyer (concurring).} The holding is narrow and does not reach broadcasts or Westboro's online ``epic''; liability here would burden public debate without proportionately advancing the State's interest.

\textbf{Alito (dissenting).} Westboro's press release, targeted signs, and online ``epic'' show a deliberate attack on a private family to win publicity. Mixing in public commentary should not immunize a personal attack, and a public street should not become a ``free-fire zone.''

\section{Examples: Future Applications}

\begin{enumerate}
\item \textbf{Same-Side (Protected): Police Funeral Protest.} Activists condemning police violence on a distant public sidewalk, following police instructions, are protected against an IIED verdict.
\item \textbf{Same-Side (Protected): Columnist and Landlord.} A savage column calling a private landlord a slumlord during a housing-safety scandal addresses a public controversy and cannot ground liability for outrageousness.
\item \textbf{Opposite-Side (Liability): Private Feud.} Yard signs mocking neighbors' miscarriage after a property-line dispute, with one tax slogan added, are a private vendetta; the claim may proceed.
\item \textbf{Opposite-Side (Likely Valid): Buffer Law.} A content-neutral ban on picketing within 300 feet of a funeral is judged as a time, place, and manner rule and likely survives, though \textit{McCullen v. Coakley} requires real tailoring.
\item \textbf{Fence-Sitter: Residential Picketing and Doxxing.} Anti-drunk-driving activists picket a grieving family's home and post their address; DUI safety is a classic public concern on a public sidewalk, but the home's heightened privacy under \textit{Frisby v. Schultz}, the doxxing campaign's resemblance to Westboro's unadjudicated ``epic,'' and possible ``true threats'' under \textit{Counterman v. Colorado} cut the other way.
\end{enumerate}

\section{Critique}

\textbf{The ``pre-existing grudge'' paradox.} As Alito observed, Roberts's reasoning means that deliberately inflicting severe distress on a grieving stranger to attract cameras receives \textit{greater} protection than lashing out in a private feud.

\textbf{Indeterminacy of ``public concern.''} Courts struggle to police the line without ad hoc content judgments. If adding a few political slogans shifts the ``dominant theme'' of targeted abuse, \textit{Snyder} may offer a blueprint for immunizing private harassment. The opinion also left open residential targeting and online campaigns, where these tensions will next arise.

\section{Key Quotations}

\begin{quote}
``Speech is powerful. It can stir people to action, move them to tears of both joy and sorrow, and---as it did here---inflict great pain. On the facts before us, we cannot react to that pain by punishing the speaker. As a Nation we have chosen a different course---to protect even hurtful speech on public issues to ensure that we do not stifle public debate.''
\end{quote}
\noindent --- Chief Justice Roberts, Opinion of the Court, 562 U.S. at 460--461.

\begin{quote}
``What Westboro said, in the whole context of how and where it chose to say it, is entitled to `special protection' under the First Amendment, and that protection cannot be overcome by a jury finding that the picketing was outrageous.''
\end{quote}
\noindent --- Chief Justice Roberts, Opinion of the Court, 562 U.S. at 458.

\begin{quote}
``Our profound national commitment to free and open debate is not a license for the vicious verbal assault that occurred in this case. \dots\ In order to have a society in which public issues can be openly and vigorously debated, it is not necessary to allow the brutalization of innocent victims like petitioner.''
\end{quote}
\noindent --- Justice Alito, Dissenting, 562 U.S. at 463, 475.

\end{document}
